\documentclass[12pt,a4paper]{article}

% --- Pacotes Básicos ---
\usepackage[utf8]{inputenc}
\usepackage[T1]{fontenc}
\usepackage[brazilian]{babel}
\usepackage{geometry}
\usepackage{setspace}
\usepackage{titlesec}
\usepackage{indentfirst}
\usepackage{helvet} % Fonte similar à Arial/Sans-Serif

% --- Configuração de Margens (Padrão Acadêmico) ---
\geometry{top=3cm, bottom=2cm, left=3cm, right=2cm}

% --- Configuração de Parágrafos e Espaçamento ---
\onehalfspacing % Espaçamento 1.5
\renewcommand{\familydefault}{\sfdefault} % Define Sans-Serif como padrão

% --- Estilo de Títulos ---
\titleformat{\section}{\normalfont\bfseries\uppercase}{\thesection}{1em}{}
\titleformat{\subsection}{\normalfont\bfseries}{\thesubsection}{1em}{}

\begin{document}

% --- CAPA ---
\begin{titlepage}
    \begin{center}
        \textbf{CENTRO UNIVERSITÁRIO – CATÓLICA DE SANTA CATARINA} \\
        \textbf{CURSO DE [NOME DO CURSO]} \\[4cm]
        
        \textbf{[NOME DOS ACADÊMICOS]} \\[4cm]
        
        \textbf{\large RELATÓRIO FINAL \\ PAC EXTENSIONISTA} \\
        \vfill
        
        \textbf{JOINVILLE, \textit{[DIA]} DE \textit{[MÊS]} DE 2025}
    \end{center}
\end{titlepage}

% --- FICHA DE APROVAÇÃO ---
\newpage
\begin{center}
    \textbf{FICHA DE APROVAÇÃO} \\[2cm]
\end{center}

\noindent RELATÓRIO do PAC Extensionista apresentado como requisito parcial de avaliação na disciplina PAC Extensionista no curso de graduação de \textit{[Nome do Curso]} do Centro Universitário Católica de Santa Catarina em Joinville, sob supervisão e orientação do professor \textit{[Nome do Professor]}.\par\vspace{2cm}

\begin{center}
    Joinville, \textit{[Dia]} de \textit{[Mês]} de 2025. \\[2cm]
    
    \rule{10cm}{0.5pt} \\
    \textit{[NOME DOS ACADÊMICOS]} \\[2cm]
    
    \rule{10cm}{0.5pt} \\
    \textit{[Nome do Professor]} \\
    Professor responsável
\end{center}

% --- SUMÁRIO ---
\newpage
\tableofcontents

% --- CONTEÚDO ---
\newpage

\section{Introdução}
% DICA: Apresentar a problemática, a justificativa e a pertinência.
% Utilize citações para fundamentar a importância do assunto.
\textit{[Escreva aqui a introdução, situando o leitor quanto ao tema, problema a ser resolvido e importância para o público beneficiado.]}

\section{Descrição do público beneficiado pelas ações de extensão}
% DICA: Descreva o perfil, local e dados do público beneficiado.
\textit{[Descreva brevemente o público, perfil e localização geográfica das ações.]}

\section{Objetivos}

\subsection{Objetivo Geral}
% DICA: Inicie com verbos de ação (Analisar, Verificar, Realizar, etc.).
\textit{[Defina a intenção maior do projeto.]}

\subsection{Objetivos Específicos}
% DICA: Máximo de 3 objetivos. Inicie com verbos de ação.
\begin{itemize}
    \item \textit{[Objetivo específico 1]};
    \item \textit{[Objetivo específico 2]};
    \item \textit{[Objetivo específico 3]}.
\end{itemize}

\section{Descrição das principais atividades realizadas}
% DICA: Relacione as atividades aos objetivos específicos (ex: 4.1, 4.2).
\textit{[Descreva o que e como as ações foram desenvolvidas para solucionar a problemática.]}

\section{Avaliação do projeto pelo público beneficiado}
% DICA: Apresente percepções, metodologia de coleta e resultados.
\textit{[Descreva a metodologia de pesquisa e os resultados obtidos com a comunidade.]}

\section{Considerações Finais -- Autoavaliação do PAC Extensionista}
% DICA: Responda se os objetivos foram alcançados, pontos fortes/fracos e aprendizados.
\textit{[Arremate final do trabalho com base nos resultados obtidos e na experiência da equipe.]}

\newpage
\section*{Referências}
\addcontentsline{toc}{section}{Referências}
% DICA: Ordem alfabética, espaçamento simples.
\textit{[Liste aqui as referências utilizadas conforme as normas da Católica SC.]}

\newpage
\section*{Apêndices}
\addcontentsline{toc}{section}{Apêndices}
\textit{[Material de autoria do pesquisador - Opcional.]}

\newpage
\section*{Anexos}
\addcontentsline{toc}{section}{Anexos}
\textit{[Material não elaborado pelo pesquisador - Opcional.]}

\end{document}